% GENERATED FILE. Edit canonical lessons, then run npm run generate:books.
% canonical-source-hash: fnv1a64:027b9bd8a4d7d3a0
% canonical-lessons: PA-C282-ataichi, PA-C282-bisatra, PA-C282-jhola, PA-C282-sambhna, PA-C282-nikalna

\chapter{Bags for the Road}
\label{ch:pa-a2-bags-for-the-road}
\hlchaptermodality{282}

\begin{chapteropening}
\textbf{By the end of this chapter:} \emph{I can talk about packing for a journey.}

Complete the last lesson of chapter 282: I can talk about packing for a journey.
\end{chapteropening}

\section[\texorpdfstring{\pa{ਅਟੈਚੀ} (aṭaichī)}{ਅਟੈਚੀ (aṭaichī)}]{\pa{ਅਟੈਚੀ} (aṭaichī) — a suitcase}

\label{lesson:PA-C282-ataichi}



\begin{quote}
Before the new one: say the Punjabi for to put off, then the Punjabi for at once.
\end{quote}

\subsection*{You'll want to know: \pa{ਅਟੈਚੀ}}
\textbf{\pa{ਅਟੈਚੀ}} — \emph{aṭaichī} — \textquotedblleft{}a suitcase\textquotedblright{}.

\begin{quote}
\pa{ਮੇਰੀ} \pa{ਅਟੈਚੀ} \pa{ਬਹੁਤ} \pa{ਭਾਰੀ} \pa{ਹੈ।}
\end{quote}

\begin{quote}
\emph{merī aṭaichī bahut bhārī hai.}
\end{quote}

\begin{quote}
My suitcase is very heavy.
\end{quote}

A feminine noun.

\begin{cousinweb}
\textbf{\pa{ਅਟੈਚੀ}} (\emph{aṭaichī}) is borrowed from English \emph{attaché}, as in an attaché case.
\end{cousinweb}

\begin{grammarlens}[title={What you've built}]
One more word for saying what comes next.
\end{grammarlens}

\subsection*{Your turn}
\begin{itemize}
\raggedright
  \item \emph{Say it:} \emph{aṭaichī}
  \item \emph{Say it:} \emph{aṭaichī}, once more
  \item \emph{Say it:} \emph{merī aṭaichī bahut bhārī hai.} — My suitcase is very heavy.
  \item \emph{Recall:} say the Punjabi for to put off, then the Punjabi for at once, then say \emph{aṭaichī} again
\end{itemize}

\begin{tcolorbox}[breakable,colback=teal!4,colframe=teal!35!black,title={Before you move on}]
What does \pa{ਅਟੈਚੀ} mean? (\textquotedblleft{}A suitcase\textquotedblright{}.) Say it once more.
\end{tcolorbox}
\section[\texorpdfstring{\pa{ਬਿਸਤਰਾ} (bisatrā)}{ਬਿਸਤਰਾ (bisatrā)}]{\pa{ਬਿਸਤਰਾ} (bisatrā) — a bedding roll; bedding}

\label{lesson:PA-C282-bisatra}



\begin{quote}
Before the new one: say the Punjabi for at once, then the Punjabi for a suitcase.
\end{quote}

\subsection*{You'll want to know: \pa{ਬਿਸਤਰਾ}}
\textbf{\pa{ਬਿਸਤਰਾ}} — \emph{bisatrā} — \textquotedblleft{}a bedding roll; bedding\textquotedblright{}.

\begin{quote}
\pa{ਅਸੀਂ} \pa{ਬਿਸਤਰਾ} \pa{ਰੇਲ} \pa{ਵਿੱਚ} \pa{ਰੱਖਾਂਗੇ।}
\end{quote}

\begin{quote}
\emph{asīṁ bisatrā rel vich rakkhāṅge.}
\end{quote}

\begin{quote}
We will put the bedding roll in the train.
\end{quote}

A masculine noun.

\begin{cousinweb}
\textbf{\pa{ਬਿਸਤਰਾ}} (\emph{bisatrā}) is a loan, usually traced to Persian \emph{bistar}, \textquotedblleft{}bedding\textquotedblright{}.
\end{cousinweb}

\begin{grammarlens}[title={What you've built}]
One more word for saying what comes next.
\end{grammarlens}

\subsection*{Your turn}
\begin{itemize}
\raggedright
  \item \emph{Say it:} \emph{bisatrā}
  \item \emph{Say it:} \emph{bisatrā}, once more
  \item \emph{Say it:} \emph{asīṁ bisatrā rel vich rakkhāṅge.} — We will put the bedding roll in the train.
  \item \emph{Recall:} say the Punjabi for at once, then the Punjabi for a suitcase, then say \emph{bisatrā} again
  \item \emph{Write it:} copy \textbf{\pa{ਬਿਸਤਰਾ}} once, checking the whole word
\end{itemize}

\begin{tcolorbox}[breakable,colback=teal!4,colframe=teal!35!black,title={Before you move on}]
What does \pa{ਬਿਸਤਰਾ} mean? (\textquotedblleft{}A bedding roll; bedding\textquotedblright{}.) Say it once more.
\end{tcolorbox}
\section[\texorpdfstring{\pa{ਝੋਲਾ} (jholā)}{ਝੋਲਾ (jholā)}]{\pa{ਝੋਲਾ} (jholā) — a cloth shoulder bag}

\label{lesson:PA-C282-jhola}



\begin{quote}
Before the new one: say the Punjabi for a suitcase, then the Punjabi for a bedding roll.
\end{quote}

\subsection*{You'll want to know: \pa{ਝੋਲਾ}}
\textbf{\pa{ਝੋਲਾ}} — \emph{jholā} — \textquotedblleft{}a cloth shoulder bag\textquotedblright{}.

\begin{quote}
\pa{ਝੋਲੇ} \pa{ਵਿੱਚ} \pa{ਪਾਣੀ} \pa{ਦੀ} \pa{ਬੋਤਲ} \pa{ਹੈ।}
\end{quote}

\begin{quote}
\emph{jhole vich pāṇī dī botal hai.}
\end{quote}

\begin{quote}
There is a bottle of water in the shoulder bag.
\end{quote}

A masculine noun.

\begin{cousinweb}
\textbf{\pa{ਝੋਲਾ}} (\emph{jholā}) is inherited from a reconstructed \emph{*jhōla}, \textquotedblleft{}bag\textquotedblright{}; CDIAL lists Punjabi \emph{jholā} under it (CDIAL 5415).
\end{cousinweb}

\begin{grammarlens}[title={What you've built}]
One more word for saying what comes next.
\end{grammarlens}

\subsection*{Your turn}
\begin{itemize}
\raggedright
  \item \emph{Say it:} \emph{jholā}
  \item \emph{Say it:} \emph{jholā}, once more
  \item \emph{Say it:} \emph{jhole vich pāṇī dī botal hai.} — There is a bottle of water in the shoulder bag.
  \item \emph{Recall:} say the Punjabi for a suitcase, then the Punjabi for a bedding roll, then say \emph{jholā} again
  \item \emph{Write it:} copy \textbf{\pa{ਝੋਲਾ}} once, checking the whole word
\end{itemize}

\begin{tcolorbox}[breakable,colback=teal!4,colframe=teal!35!black,title={Before you move on}]
What does \pa{ਝੋਲਾ} mean? (\textquotedblleft{}A cloth shoulder bag\textquotedblright{}.) Say it once more.
\end{tcolorbox}
\section[\texorpdfstring{\pa{ਸਾਂਭਣਾ} (sāmbhṇā)}{ਸਾਂਭਣਾ (sāmbhṇā)}]{\pa{ਸਾਂਭਣਾ} (sāmbhṇā) — to keep safe; to look after}

\label{lesson:PA-C282-sambhna}



\begin{quote}
Before the new one: say the Punjabi for a bedding roll, then the Punjabi for a cloth shoulder bag.
\end{quote}

\subsection*{You'll want to know: \pa{ਸਾਂਭਣਾ}}
\textbf{\pa{ਸਾਂਭਣਾ}} — \emph{sāmbhṇā} — \textquotedblleft{}to keep safe; to look after\textquotedblright{}.

\begin{quote}
\pa{ਮੈਂ} \pa{ਟਿਕਟਾਂ} \pa{ਸਾਂਭਾਂਗਾ।}
\end{quote}

\begin{quote}
\emph{maiṁ ṭikṭāṁ sāmbhāṅgā.}
\end{quote}

\begin{quote}
I will keep the tickets safe.
\end{quote}

Past (I …): \pa{ਮੈਂ} \pa{ਟਿਕਟ} \pa{ਸਾਂਭੀ}, \emph{maiṁ ṭikaṭ sāmbhī}.

\begin{cousinweb}
\textbf{\pa{ਸਾਂਭਣਾ}} (\emph{sāmbhṇā}) is inherited from Sanskrit \emph{saṁbhṛta}, \textquotedblleft{}collected\textquotedblright{}; CDIAL lists Punjabi \emph{sāmbhṇā}, \textquotedblleft{}to support, take care of\textquotedblright{}, under it (CDIAL 12970).
\end{cousinweb}

\begin{grammarlens}[title={What you've built}]
One more word for saying what comes next.
\end{grammarlens}

\subsection*{Your turn}
\begin{itemize}
\raggedright
  \item \emph{Say it:} \emph{sāmbhṇā}
  \item \emph{Say it:} \emph{sāmbhṇā}, once more
  \item \emph{Say it:} \emph{maiṁ ṭikṭāṁ sāmbhāṅgā.} — I will keep the tickets safe.
  \item \emph{Recall:} say the Punjabi for a bedding roll, then the Punjabi for a cloth shoulder bag, then say \emph{sāmbhṇā} again
  \item \emph{Write it:} copy \textbf{\pa{ਸਾਂਭਣਾ}} once, checking the whole word
\end{itemize}

\begin{tcolorbox}[breakable,colback=teal!4,colframe=teal!35!black,title={Before you move on}]
What does \pa{ਸਾਂਭਣਾ} mean? (\textquotedblleft{}To keep safe; to look after\textquotedblright{}.) Say it once more.
\end{tcolorbox}
\section[\texorpdfstring{\pa{ਨਿਕਲਣਾ} (nikalṇā)}{ਨਿਕਲਣਾ (nikalṇā)}]{\pa{ਨਿਕਲਣਾ} (nikalṇā) — to come out; to head out}

\label{lesson:PA-C282-nikalna}



\begin{quote}
Before the new one: say the Punjabi for a cloth shoulder bag, then the Punjabi for to keep safe.
\end{quote}

\subsection*{You'll want to know: \pa{ਨਿਕਲਣਾ}}
\textbf{\pa{ਨਿਕਲਣਾ}} — \emph{nikalṇā} — \textquotedblleft{}to come out; to head out\textquotedblright{}.

\begin{quote}
\pa{ਅਸੀਂ} \pa{ਸਵੇਰੇ} \pa{ਛੇ} \pa{ਵਜੇ} \pa{ਘਰ} \pa{ਤੋਂ} \pa{ਨਿਕਲਾਂਗੇ।}
\end{quote}

\begin{quote}
\emph{asīṁ savere chhe vaje ghar ton niklāṅge.}
\end{quote}

\begin{quote}
We will head out from home at six in the morning.
\end{quote}

Past (I …): \pa{ਮੈਂ} \pa{ਘਰ} \pa{ਤੋਂ} \pa{ਨਿਕਲਿਆ}, \emph{maiṁ ghar ton nikliā}.

\begin{cousinweb}
\textbf{\pa{ਨਿਕਲਣਾ}} (\emph{nikalṇā}) is inherited from a reconstructed \emph{*niṣkalati}, \textquotedblleft{}comes out\textquotedblright{}; CDIAL lists Punjabi \emph{nikkalṇā} under it (CDIAL 7478).
\end{cousinweb}

\begin{grammarlens}[title={What you've built}]
One more word for saying what comes next.
\end{grammarlens}

\subsection*{Your turn}
\begin{itemize}
\raggedright
  \item \emph{Say it:} \emph{nikalṇā}
  \item \emph{Say it:} \emph{nikalṇā}, once more
  \item \emph{Say it:} \emph{asīṁ savere chhe vaje ghar ton niklāṅge.} — We will head out from home at six in the morning.
  \item \emph{Recall:} say the Punjabi for a cloth shoulder bag, then the Punjabi for to keep safe, then say \emph{nikalṇā} again
  \item \emph{Write it:} copy \textbf{\pa{ਨਿਕਲਣਾ}} once, checking the whole word
\end{itemize}

\begin{tcolorbox}[breakable,colback=teal!4,colframe=teal!35!black,title={Before you move on}]
What does \pa{ਨਿਕਲਣਾ} mean? (\textquotedblleft{}To come out; to head out\textquotedblright{}.) Say it once more.
\end{tcolorbox}
